\documentclass[11pt, a4paper]{article}
\usepackage[utf8]{inputenc}
\usepackage[spanish]{babel}
\usepackage[margin=2cm]{geometry}
\usepackage{tabularx}
\usepackage{amssymb}
\usepackage[most]{tcolorbox}

\begin{document}
\begin{center}
    {\Large \textbf{Formulario de Diagnóstico y Punto de Control Post-Hito 1}}\\[0.2cm]
    {\large Circuitos Electrónicos III (IMT-221) - Cierre de Unidad 1}
\end{center}
\hrule
\vspace{0.4cm}

\begin{tcolorbox}[colback=white, colframe=black, title=\textbf{Resumen de la Sesión}]
\textbf{Fortalezas observadas en la clase:}\\[0.3cm] \hrulefill \\[0.3cm]
\textbf{Debilidades generalizadas de la clase:}\\[0.3cm] \hrulefill \\[0.3cm]
\textbf{Estudiantes/Grupos que requieren recuperación inmediata:}\\[0.3cm] \hrulefill 
\end{tcolorbox}

\begin{tcolorbox}[colback=white, colframe=black, title=\textbf{Estado por Dominio Técnico}]
\begin{tabularx}{\textwidth}{l X X X X}
\textbf{Protección Zener / Clamping:} & $\square$ Logrado & $\square$ En progreso & $\square$ Débil & $\square$ Desconocido \\
\textbf{Transitorios Inductivos:} & $\square$ Logrado & $\square$ En progreso & $\square$ Débil & $\square$ Desconocido \\
\textbf{Diodo Flyback:} & $\square$ Logrado & $\square$ En progreso & $\square$ Débil & $\square$ Desconocido \\
\textbf{Competencia Analítica / Python:} & $\square$ Logrado & $\square$ En progreso & $\square$ Débil & $\square$ Desconocido \\
\textbf{Competencia LTspice:} & $\square$ Logrado & $\square$ En progreso & $\square$ Débil & $\square$ Desconocido \\
\textbf{Implementación de Hardware:} & $\square$ Logrado & $\square$ En progreso & $\square$ Débil & $\square$ Desconocido \\
\textbf{Validación Experimental / CSV:} & $\square$ Logrado & $\square$ En progreso & $\square$ Débil & $\square$ Desconocido \\
\end{tabularx}
\end{tcolorbox}

\begin{tcolorbox}[colback=white, colframe=black, title=\textbf{Registro Cualitativo}]
\textbf{Brechas de evidencia observadas (Datos faltantes, reportes incompletos):}\\[0.4cm] \hrulefill \\[0.4cm]
\textbf{Conceptos erróneos recurrentes:}\\[0.4cm] \hrulefill \\[0.4cm]
\textbf{Problemas logísticos, de componentes o instrumentos (ej. osciloscopio):}\\[0.4cm] \hrulefill 
\end{tcolorbox}

\begin{tcolorbox}[colback=blue!5, colframe=blue!75!black, title=\textbf{Decisión Pedagógica para la Siguiente Sesión}]
$\square$ Transición directa a la estructura del transistor BJT\\
$\square$ Refuerzo delimitado de Unidad 1 (10–20 min) + Inicio de BJT\\
$\square$ Remedición sustantiva requerida antes de introducir BJT\\[0.3cm]
\textbf{Contenido mínimo de remedición (si aplica):}\\[0.3cm] \hrulefill \\[0.3cm]
\textbf{¿Es factible el Hito 2 en la Semana 8 bajo el ritmo actual?}\\
$\square$ Sí \hspace{1cm} $\square$ En riesgo \hspace{1cm} $\square$ Requiere ajuste de cronograma
\end{tcolorbox}

\end{document}
